\chapter{Python の練習問題の解答例}
\label{app:python-answers}

\ref{chap:python_basics}章と\ref{chap:python_class}章の練習問題の解答例です。
解答例は一例です。
同じ結果になれば、違う書き方でも構いません。
解答例を見る前に、まずは自分で考えて、プログラムを書いて動かしてみましょう。

\section{第18章 Python の基本}
\label{sec:python-basics-answers}

\begin{description}
  \item[練習問題 18.1 (1)] たとえば \cmd{5 + 8 - 6 / 2} です（$6 / 2 = 3$ が先に計算されるので $5 + 8 - 3 = 10$）。\cmd{/} の計算結果は小数になるので、Python で実行すると \cmd{10.0} と表示されます。
  \item[練習問題 18.1 (2)] \cmd{200 // 60} で 3（分）、\cmd{200 \% 60} で 20（秒）なので、3 分 20 秒です。
\end{description}

\begin{lstlisting}[style=python, caption={練習問題 18.2 (1) の解答例}, label={lst:python-basics-a3}]
money = 4000
book = 980
chips = 250
ice = 150
book_tax = 0.10
food_tax = 0.08

total = book * 2 * (1 + book_tax) + (chips + ice) * (1 + food_tax)
print(round(money - total))   # 1412
\end{lstlisting}

\begin{lstlisting}[style=python, caption={練習問題 18.2 (2)・(3) の解答例}, label={lst:python-basics-a4}]
# 練習問題 18.2 (2)
title = input("アニメ名=> ")
print(f"{title}は{len(title)}文字です。")

# 練習問題 18.2 (3)
radius = float(input("車輪の半径 [m] => "))
rps = float(input("1 秒あたりの回転数 => "))
print(f"速さは {2 * 3.14 * radius * rps:.2f} m/s です")
\end{lstlisting}

\begin{lstlisting}[style=python, caption={練習問題 18.3 (1)〜(4) の解答例}, label={lst:python-basics-a6}]
# 練習問題 18.3 (1)
colors = ["green", "blue", "red", "black", "white"]
print(colors[1:4])

# 練習問題 18.3 (2)
members = ["satou", "tanaka", "suzuki"]
members.append("itou")
members.remove("satou")
print(members)       # ['tanaka', 'suzuki', 'itou']
print(members[1])    # suzuki

# 練習問題 18.3 (3)
heights = [176, 158, 164]
print(sum(heights) / len(heights))   # 166.0

# 練習問題 18.3 (4)
robot = {"name": "turtle", "max_speed": 0.5}
robot["max_speed"] = 0.8
robot["battery"] = 12.0
print(f"{robot['name']} の最高速度は {robot['max_speed']} m/s")
\end{lstlisting}

f 文字列の \cmd{\{ \}} の中で辞書のキーを書くときは、外側の \cmd{"} とぶつからないように \cmd{'} を使います。

\begin{lstlisting}[style=python, caption={練習問題 18.4 (1)〜(4) の解答例}, label={lst:python-basics-a10}]
# 練習問題 18.4 (1)
a = float(input("a => "))
b = float(input("b => "))
if a > b:
    print("a の方が大きい")
elif a < b:
    print("b の方が大きい")
else:
    print("a と b は同じ")

# 練習問題 18.4 (2)
height = float(input("身長 [m] => "))
weight = float(input("体重 [kg] => "))
bmi = weight / height ** 2
print(f"BMI = {bmi:.1f}")
if bmi < 18.5:
    print("やせ")
elif bmi < 25:
    print("標準")
elif bmi < 30:
    print("肥満")
else:
    print("高度肥満")

# 練習問題 18.4 (3)
num = int(input("整数 => "))
if num % 2 == 0:
    print("偶数である")
    if num >= 25:
        print("25 以上である")
if num % 3 == 0:
    print("3 の倍数である")

# 練習問題 18.4 (4)
voltage = float(input("電圧 [V] => "))
if voltage <= 0 or voltage > 15:
    print("センサの値が異常です")
elif voltage >= 12.0:
    print("充電十分")
elif voltage >= 11.0:
    print("そろそろ充電")
else:
    print("すぐに充電してください")
\end{lstlisting}

練習問題 18.4 (3) のように、if 文の中に if 文を書くこともできます（インデントを 1 段深くします）。

\begin{lstlisting}[style=python, caption={練習問題 18.5 (1)〜(4) の解答例}, label={lst:python-basics-a14}]
# 練習問題 18.5 (1)
total = 0
for i in range(1, 31):
    total += i
print(total)   # 465

# 練習問題 18.5 (2)
nums = [1, 62, 3, 6743, 51, 6783, 42, 4671, 5, 236]
for n in nums:
    if n >= 1000:
        continue
    print(n)

# 練習問題 18.5 (3)
n = int(input("数> "))
for i in range(1, n + 1):
    print(f"{i}:" + "■" * i)

# 練習問題 18.5 (4)
count = 0
total = 0
while True:
    data = float(input("データ入力（負の数で終了）> "))
    if data < 0:
        break
    count += 1
    total += data
if count > 0:
    print(f"個数: {count} 合計: {total} 平均: {total / count}")
\end{lstlisting}

\begin{lstlisting}[style=python, caption={練習問題 18.5 (5)・(6) の解答例}, label={lst:python-basics-a18}]
import random

# 練習問題 18.5 (5)
answer = random.randint(1, 30)
for i in range(5):
    guess = int(input("値を入力してください=> "))
    if guess == answer:
        print("正解です！")
        break
    elif guess < answer:
        print("残念！それよりも大きいです！")
    else:
        print("残念！それよりも小さいです！")
print(f"正解は {answer} でした。")

# 練習問題 18.5 (6)
hands = ["グー", "チョキ", "パー"]
results = ["あいこ", "負け", "勝ち"]
mine = int(input("あなたの手は？（0:グー, 1:チョキ, 2:パー）=> "))
computer = random.randint(0, 2)
print(f"あなた: {hands[mine]}  コンピュータ: {hands[computer]}")
print(results[(mine - computer + 3) % 3])
\end{lstlisting}

練習問題 18.5 (6) では、手の名前や結果をリストに入れておき、インデックスで取り出しています。
if 文を何重にも書くよりも、短くわかりやすいプログラムになります。

\begin{lstlisting}[style=python, caption={練習問題 18.6 (1)〜(5) の解答例}, label={lst:python-basics-a20}]
# 練習問題 18.6 (1)
def motto():
    print("完璧を")
    print("目指すよりも")
    print("まずは")
    print("終わらせろ")


motto()
motto()
motto()


# 練習問題 18.6 (2)
def even_or_odd(num):
    if num % 2 == 0:
        print("偶数")
    else:
        print("奇数")


even_or_odd(5)
even_or_odd(20)
even_or_odd(1)


# 練習問題 18.6 (3)
def maximum(a, b, c):
    largest = a
    if b > largest:
        largest = b
    if c > largest:
        largest = c
    return largest


print(maximum(40, 23, 60))   # 60


# 練習問題 18.6 (4)
def find_max_min(values):
    return max(values), min(values)


largest, smallest = find_max_min([3, 62, 4, 56, 346, 8, 2, 893, 52, 46])
print(largest, smallest)   # 893 2


# 練習問題 18.6 (5)
def valid_average(distances):
    valid = []
    for d in distances:
        if d >= 0:
            valid.append(d)
    if len(valid) == 0:
        return None
    return sum(valid) / len(valid)


print(valid_average([1.2, -1.0, 0.8]))   # 1.0
print(valid_average([-1.0, -1.0]))       # None
\end{lstlisting}

\section{第19章 クラスとモジュール}
\label{sec:python-class-answers}

練習問題 19.7（最終課題）には決まった答えがないので、解答例は載せていません。

\begin{lstlisting}[style=python, caption={練習問題 19.1 (1)・(2) の解答例}, label={lst:python-class-a1}]
# 練習問題 19.1 (1)
class Human:
    def __init__(self, name, age):
        self.name = name
        self.age = age

    def introduce(self):
        print(f"私の名前は{self.name}です")

    def show_age(self):
        print(f"年齢は{self.age}歳です")


john = Human("ジョン", 23)
marine = Human("マリン", 18)
john.introduce()
john.show_age()
marine.introduce()
marine.show_age()


# 練習問題 19.1 (2)
class Tax:
    def __init__(self, rate):
        self.rate = rate   # 税率 [%]

    def price(self, value):
        total = round(value * (1 + self.rate / 100))
        print(f"{value} 円（税率 {self.rate}%）→ {total} 円")


tax8 = Tax(8)
tax10 = Tax(10)
for value in [250, 580]:
    tax8.price(value)
    tax10.price(value)
\end{lstlisting}

\begin{lstlisting}[style=python, caption={練習問題 19.1 (3) の解答例}, label={lst:python-class-a3}]
class Robot:
    def __init__(self, name, max_speed=0.5):
        self.name = name
        self.max_speed = max_speed
        self.x = 0.0
        self.battery = 100

    def move(self, speed, duration):
        if self.battery <= 0:
            print("バッテリ切れです")
            return
        if speed > self.max_speed:
            speed = self.max_speed
        distance = speed * duration
        self.x += distance
        self.battery -= distance * 10
        print(f"x = {self.x:.2f} m, バッテリ {self.battery:.0f}%")


robot = Robot("turtle")
for i in range(12):
    robot.move(0.5, 2.0)
\end{lstlisting}

\cmd{return} だけを書くと、戻り値を返さずに、その時点で関数（メソッド）を終了します。

\begin{lstlisting}[style=python, caption={練習問題 19.2 (1)・(2) の解答例}, label={lst:python-class-a4}]
import random


class Player:
    def __init__(self, name, hp):
        self.name = name
        self.hp = hp

    def attack(self):
        print(f"{self.name}は攻撃をした！")


# 練習問題 19.2 (1)
class Wizard(Player):
    def __init__(self, name, hp, mp):
        super().__init__(name, hp)
        self.mp = mp

    def magic(self):
        if self.mp < 10:
            print("魔力が足りない")
            return
        self.mp -= 10
        print(f"{self.name}は魔法を唱えた！（残り魔力: {self.mp}）")


wizard = Wizard("魔法使い", 60, 25)
wizard.attack()
wizard.magic()
wizard.magic()
wizard.magic()


# 練習問題 19.2 (2)
class Sensor:
    def __init__(self, name):
        self.name = name

    def read(self):
        return 0.0


class DummyLidar(Sensor):
    def read(self):
        return random.uniform(0.1, 5.0)


lidar = DummyLidar("lidar")
print(f"{lidar.name}: {lidar.read():.2f} m")
\end{lstlisting}

\cmd{DummyLidar} では \cmd{\_\_init\_\_} を定義していないので、親クラスの \cmd{\_\_init\_\_} がそのまま使われます。

\begin{lstlisting}[style=python, caption={練習問題 19.3 (1)・(2) の解答例}, label={lst:python-class-a6}]
# 練習問題 19.3 (1)
def apply_all(values, func):
    results = []
    for v in values:
        results.append(func(v))
    return results


def to_cm(m):
    return m * 100


print(apply_all([1.0, 2.5, 4.0], to_cm))   # [100.0, 250.0, 400.0]


# 練習問題 19.3 (2)
def watch(distances, on_obstacle):
    for d in distances:
        if d < 0.3:
            on_obstacle(d)


def warn(d):
    print(f"障害物を検出: {d} m")


watch([1.2, 0.25, 0.8, 0.1], warn)
\end{lstlisting}

\begin{lstlisting}[style=python, caption={練習問題 19.4 (1)・(2) の解答例}, label={lst:python-class-a8}]
import math
from math import pi

# 練習問題 19.4 (1)
print(pi * 10 * 10)                # 314.1592653589793

# 練習問題 19.4 (2)
print(math.radians(30 * 4))        # 2.0943951023931953
\end{lstlisting}

\begin{lstlisting}[style=python, caption={練習問題 19.4 (3) の解答例（geometry.py）}, label={lst:python-class-a10}]
import math


def distance(x1, y1, x2, y2):
    """2 点間の距離を返す"""
    return math.sqrt((x2 - x1) ** 2 + (y2 - y1) ** 2)


def normalize_angle(deg):
    """角度を -180 以上 180 未満に正規化する"""
    return (deg + 180) % 360 - 180
\end{lstlisting}

\begin{lstlisting}[style=python, caption={練習問題 19.4 (3) の解答例（geometry を使うファイル）}, label={lst:python-class-a10b}]
import geometry

print(geometry.distance(0, 0, 3, 4))    # 5.0
print(geometry.normalize_angle(270))    # -90
\end{lstlisting}

\begin{lstlisting}[style=python, caption={練習問題 19.5 (1)・(2) の解答例}, label={lst:python-class-a11}]
import random

# 練習問題 19.5 (1)
for i in range(10):
    num = random.randint(0, 3)
    try:
        print(100 / num)
    except ZeroDivisionError:
        print("0 で割ろうとしました")

# 練習問題 19.5 (2)
try:
    num = int(input("整数を入力してください=> "))
except ValueError:
    num = 10
print(num + 100)
\end{lstlisting}

\begin{description}
  \item[練習問題 19.6 (1)] \cmd{random.randint(1, 10)} は 10 も含むので \cmd{random.randint(1, 9)} に直します。また、ループの中の \cmd{print} で 5 もすでに表示されているので、最後の \cmd{print} を消す（またはループの中の \cmd{print} を \cmd{if} の後ろに移す）と、5 が 2 回表示されなくなります。
  \item[練習問題 19.6 (2)] \cmd{input( )} の値は文字列なので \cmd{range(int(num))} にする（1 か所目）、条件を \cmd{i \% 5 == 2 or i \% 5 == 3} にする（2 か所目、\ref{sec:python-basics-if}節の注意を参照）、\cmd{print("i")} は文字 \cmd{i} を表示してしまうので \cmd{print(i)} にする（3 か所目）。
  \item[練習問題 19.6 (3)] 次の 5 か所を直します。
    \begin{enumerate}
      \item \cmd{input( )} の値を \cmd{int( )} で整数に変換する
      \item 条件の \cmd{and} を \cmd{or} にする（\cmd{player} が 0 かつ 1 かつ 2 になることはない）
      \item \cmd{prin} を \cmd{print} にする
      \item \cmd{random.randint(0, 3)} を \cmd{random.randint(0, 2)} にする（3 だと \cmd{hands[3]} で \cmd{IndexError} になる）
      \item \cmd{dceision} を \cmd{decision} にする
    \end{enumerate}
\end{description}
